\section{Introduction}
\label{sec:intro}

What happens to exchange rates when a country imposes tariffs? The conventional answer, going back to \citet{lerner1936}, \citet{metzler1949}, \citet{mundell1961} and \citet{krugman1982}, is that the tariffing country's currency appreciates. A tariff shifts expenditure away from imports and improves the trade balance at the prevailing exchange rate. Under the Marshall--Lerner condition, restoring external balance then requires an appreciation. The prediction is robust \emph{within} the two-country model: it follows from monotonicity properties of import and export demand rather than from functional forms.

The 2025 tariff episode is hard to square with this prediction. After the April~2 announcements the dollar weakened, notably against the euro, while moving little against the renminbi. Much of the recent literature has turned to financial explanations. These include tariff uncertainty and currency risk premia \citep{kalemliozcan2026}, the erosion of the dollar's convenience yield and hedging value \citep{acharya2026,jiang2025}, and valuation effects on dollar liabilities \citep{itskhoki2025}. \citet{ostry2025} document a striking regularity across the 2018--20 and 2025 episodes: the dollar appreciates when US tariffs are unilateral and depreciates when partners retaliate, and in US--China tit-for-tat episodes it is roughly stable against the renminbi. They interpret this pattern through dominant-currency pricing.

This paper asks a prior question: \textbf{what does the real trade channel itself imply for exchange rates once trade is multilateral?} The two-country model cannot answer it, because it has only one exchange rate, and substitution away from a tariffed supplier can only go toward domestic goods. With more than two countries, three things that coincide when $N=2$ come apart: the adjustment of trade balances to exchange rates, the allocation of tariff-induced expenditure switching across foreign suppliers, and the distinction between bilateral and effective exchange rates. Our main message is that \textbf{the conventional wisdom survives multilateralization as a statement about the effective exchange rate but fails as a statement about bilateral exchange rates.} The model characterizes precisely the substitution, network and policy conditions under which each holds.

We work in a static, long-run, real environment: one factor, linear technologies, balanced trade and no nominal rigidities. In this class of models the nominal exchange rate equals the relative wage in a common currency, and therefore the bilateral terms of trade. This is deliberate. It isolates what trade flows alone pin down, and it lets us obtain sharp analytical results. Our results should be read as the long-run exchange-rate position that the trade channel requires once prices have adjusted, conditional on the tariffs persisting.

\paragraph{Results.} We proceed in three steps.

\emph{General framework (\cref{sec:framework}).} In any differentiable $N$-country trade environment, the first-order exchange-rate response to a tariff change is $d\nu=(-J)^{-1}F\,ds$. Here $F$ is the vector of trade-balance disturbances generated by the tariff at unchanged exchange rates, and $J$ is the Jacobian of trade balances with respect to exchange rates. We define a \emph{multilateral Marshall--Lerner} (MML) condition: $-J$ is a nonsingular M-matrix. It is implied by gross substitutability of trade balances in currency values, and it is the natural heir of the Hicks--Metzler--Mosak stability conditions. MML ensures that exchange-rate adjustment is stabilizing. It does not determine the sign of any bilateral response, because $F$ is now a vector whose entries can have different signs.

\emph{The symmetric case (\cref{sec:symmetric}).} When the tariffing country's partners are symmetric, any tariff vector splits into \emph{average protection}, the common component, and \emph{relative treatment}, the zero-sum component. These are the eigen-directions of both $J$ and the tariff-incidence matrix. Average protection alone determines the nominal effective exchange rate (NEER); relative treatment alone determines the dispersion of bilateral rates around it. In a nested-CES model with a home--import elasticity $\eta$ and a cross-origin elasticity $\rho$, three margins govern everything:
\begin{itemize}
\item the \emph{domestic diversion margin} $\rho_D=\alpha_D\eta+(1-\alpha_D)(1-\alpha_T)$, which combines substitution toward domestic tradables with the income effect of rebated tariff revenue;
\item the cross-origin elasticity $\rho$;
\item a \emph{multilateral adjustment margin} $M_N$.
\end{itemize}
A unilateral tariff on one partner always appreciates the NEER and always appreciates the currency against the target. It depreciates the currency against each untariffed bystander if and only if $\rho>N\rho_D$. With realistic home bias this requires $\rho>\eta$, so the reversal is impossible in any single-elasticity specification. Bilateral retaliation lowers the bystander threshold from $N\rho_D$ to $\rho_D$ and leaves the belligerents' bilateral rate unchanged; both belligerents then depreciate against every bystander and in effective terms. This is exactly the configuration-dependence that \citet{ostry2025} document, here produced with flexible prices and no financial channel.

Finally, we show that the symmetric $N$-country economy reduces \emph{exactly} to a three-country economy: the tariffing country, the target, and a composite bystander with weight $(N-2)/(N-1)$ in the tariffing country's imports and an effective home--import elasticity that blends $\eta$ and $\rho$. This three-country representation is our workhorse for intuition. It shows that $N$ raises the bystander threshold not because the bystander bloc is larger, but because the bloc's trade is increasingly internal and highly substitutable. Bystander size alone moves the threshold in the opposite direction.

\emph{Asymmetry (\cref{sec:asymmetry}).} Away from symmetry the decomposition no longer separates cleanly, and some intuitive conclusions of the symmetric case can fail. We focus on the aggregate terms of trade. At free trade it equals the first-order welfare effect of a tariff, and under symmetry a unilateral tariff always improves it. With asymmetric partners it can deteriorate. A tariff on a partner that is small, or carries little weight in the tariffing country's imports, can produce a bystander depreciation that is large and heavily weighted enough to outweigh the gain against the target when $\rho$ is high. The trade-weighted NEER can likewise depreciate, and the relationship with $\rho$ need not even be monotone. We then revisit the 2025 episode in a deliberately stylized way, without calibrating.

\paragraph{What we do not do.} The model abstracts from nominal rigidities and invoicing, from capital, intermediates and supply chains, from trade imbalances and asset trade, and from risk and policy uncertainty. We do not attempt to fit the 2025 data. The point of the paper is conceptual: one cannot appeal to the two-country prediction to conclude that an observed multilateral depreciation must reflect financial or non-trade forces. Conversely, the framework says when such an attribution \emph{is} warranted.

\subsection{Related literature}
\label{sec:literature}

A rapidly growing literature studies the exchange-rate effects of the 2025 tariffs. One strand emphasizes financial channels: tariff uncertainty and risk premia \citep{kalemliozcan2026}, the erosion of Treasury convenience yields \citep{acharya2026,jiang2025}, and valuation effects on dollar liabilities \citep{itskhoki2025}. A second strand documents that the dollar's response depends on retaliation \citep{ostry2025}, interpreted through dominant-currency pricing. A third quantifies the welfare and third-country effects of tariffs in rich quantitative models \citep{ignatenko2025,darracqparies2026,fajgelbaum2026}. What is missing is a characterization of what the real trade channel \emph{itself} implies for bilateral and effective exchange rates once trade is multilateral. We provide that characterization. It shows that the configuration-dependence documented by \citet{ostry2025} arises in a frictionless real model, and it identifies when observed movements, such as the late-2025 dollar weakness under nearly uniform tariffs, cannot be attributed to trade.

\paragraph{Tariffs and exchange rates.} The classical two-country analysis of tariffs under flexible exchange rates \citep{lerner1936,metzler1949,mundell1961,dornbusch1974,eichengreen1981,krugman1982} delivers the appreciation result; \citet{dornbusch1974} is particularly close to us in giving non-traded goods a role in the incidence of a tariff, which enters our domestic diversion margin through the rebate. \citet{costinot2019} give the modern statement of Lerner symmetry; since our results concern the real allocation, they are unaffected by Lerner-type reformulations of the policy. A large open-macroeconomy literature studies tariffs in the short run with nominal rigidities, where monetary policy, invoicing and the fiscal use of revenue shape the response \citep{linde2019,barattieri2021,erceg2023,jeanne2021,bergin2023,auray2025,bianchi2025,werning2025}. We complement that work with the long-run real benchmark in a multilateral setting, the position toward which those short-run dynamics converge once prices have adjusted.

\paragraph{The 2025 dollar.} The financial and uncertainty explanations of the 2025 dollar \citep{kalemliozcan2026,acharya2026,itskhoki2025,cardani2025,jiang2025} are complements to our analysis, not competitors; \citet{kalemliozcan2026}, for instance, work in a two-country model, where tariff levels appreciate the currency and tariff volatility depreciates it. \citet{itskhoki2025} are the closest theory paper on the dollar: in their static benchmark a tariff appreciates the currency, and with dollar-denominated US liabilities that appreciation carries a valuation cost that sharply lowers the optimal macro tariff; a tariff war can then cost the United States its exorbitant privilege. Their mechanism is valuation and convenience yield; ours is real trade diversion. The paper closest to ours empirically is \citet{ostry2025}. Building a database of tariff announcements and retaliations for 2018--20 and 2025, they show that the dollar appreciates on unilateral US tariffs and depreciates when partners retaliate, while in US--China tit-for-tat exchanges it is roughly flat bilaterally. They interpret the pattern through dominant-currency pricing \citep{gopinath2020}. Our bilateral-retaliation result predicts exactly the same pattern, with $d\nu_B=0$ against the retaliating partner and depreciation against bystanders whenever $\rho>\rho_D$, without any nominal friction. The two explanations can coexist. Ours shows that the configuration-dependence they document is what the real trade channel requires in the long run, so it need not be read as evidence of a nominal or financial friction, and it identifies the substitution conditions under which that reading is warranted.

\paragraph{Third-country effects.} Because the exchange rate \emph{is} the terms of trade in our model class, the preferential-trade-agreement literature on third-country terms-of-trade effects applies directly \citep{viner1950,mundell1964,kemp1976,krugman1991,bond1996,bagwell1999,bagwell2002,yi2000,chang2002,ornelas2005}. Our contribution to it is narrow: the competing-exporter result, in which the excluded exporter's terms of trade always improve, becomes conditional ($\rho>N\rho_D$) once the importer also produces a competing good. Quantitative and empirical work on the third-country effects of trade wars \citep{darracqparies2026,ignatenko2025,fajgelbaum2024,freund2024,alfaro2023,fajgelbaum2026} quantifies these effects in rich environments; \citet{darracqparies2026} conclude that tariff incidence is intrinsically multilateral and cannot be inferred from two-country analysis. These papers quantify; we characterize signs.

\paragraph{Elasticities.} Our thresholds turn on the gap between the home--import (macro) and cross-origin (micro) elasticities. The evidence that cross-foreign elasticities exceed home--foreign ones \citep{broda2006,feenstra2018,li2022}, and that long-run trade elasticities are larger than short-run ones \citep{boehm2023}, is exactly the distinction our results require.

\paragraph{Multi-market stability.} The MML condition descends from the Hicks--Metzler--Mosak conditions for the stability of multiple markets and from the gross-substitutes stability results \citep{hicks1939,mosak1944,metzler1945,arrow1958,arrow1959}. We apply that lineage to the tariff--exchange-rate mapping rather than inventing the condition; the IMF's multilateral exchange rate model \citep{artus1973} is an early institutional antecedent of multilateral exchange-rate modelling of this kind.

\paragraph{Roadmap.} \Cref{sec:framework} develops the general framework. \Cref{sec:symmetric} analyzes the symmetric case and its three-country representation. \Cref{sec:asymmetry} turns to asymmetry and the 2025 episode, and \cref{sec:conclusion} concludes. \Cref{app:linear} collects the linearization, \cref{app:proofs} the proofs, and \cref{app:departures} departures from the symmetric free-trade benchmark. Online \cref{app:online} contains further robustness results.
